\documentclass[11pt]{article}
\usepackage[margin=1in]{geometry}
\usepackage{amsmath,amssymb,amsthm,mathtools}
\usepackage{graphicx}
\usepackage{microtype}
\usepackage[hidelinks]{hyperref}

\newtheorem{theorem}{Theorem}
\newtheorem{lemma}{Lemma}
\newtheorem{remark}{Remark}
\newcommand{\N}{\mathbb N}
\newcommand{\K}{\mathbb K}
\newcommand{\ellp}{\ell_p}
\newcommand{\spn}{\operatorname{span}}

\title{Devaney-chaotic Rolewicz sums are\\densely distributionally chaotic}
\author{Candidate full solution to the conjecture in arXiv:1504.02445}
\date{12 August 2026}

\begin{document}
\maketitle

\noindent\textbf{Status.} Candidate full solution, pending expert review.

\medskip
\noindent\textbf{Source.} D.~Bongiorno, U.~B. Darji, and L.~Di Piazza,
\emph{Rolewicz-type chaotic operators}~\cite{BDD}.  The conjecture is on
printed page~4 of the source PDF (arXiv:1504.02445).

\section{The question}
For a strictly increasing map $f:\N\to\N$, define
\[
 T_f(x_1,x_2,\ldots)=(x_{f(1)},x_{f(2)},\ldots)
 \qquad (x\in\ell_p).
\]
Two such maps $f,g$ are \emph{almost disjoint} in the source's sense if
$\{k:f(k)=g(k)\}$ is finite and $f(k)\ne g(k')$ whenever $k\ne k'$.
The source conjectures that, for pairwise almost-disjoint
$f_1,\ldots,f_t$ and real $c_1,\ldots,c_t$, Devaney chaos of
\[
             T=\sum_{i=1}^t c_iT_{f_i}:\ell_p\longrightarrow\ell_p
\tag{1}\label{eq:T}
\]
implies distributional chaos.

\begin{figure}[ht]
\centering
\includegraphics[width=\textwidth]{figures/open_problem_crop.png}
\caption{The conjecture in the original paper, printed page~4.  This is a
full-width crop rendered from the archived source PDF.}
\end{figure}

\paragraph{Transcription.}
``We conjecture that if $f_1,\ldots,f_t$ are strictly increasing pairwise
almost disjoint functions and $c_1,\ldots,c_t$ are real numbers, then
$T=\sum_{i=1}^t c_iT_{f_i}$ being chaotic (in the sense of Devaney) implies
that $T$ is distributionally chaotic.''

\section{Proof intuition}
Far enough along the coordinates, almost disjointness turns $T$ into a
perfect branching backward shift: every input coordinate is used in at most
one branch and row.  The norm of a tail row is therefore exactly the dual
norm
\[
 C=\|(c_1,\ldots,c_t)\|_{p'}.
\]
Hypercyclicity passes to this tail factor and forces $C>1$.

There may be collisions in finitely many initial rows.  Hypercyclicity forces
each such row to have a nonzero coefficient that eventually escapes into the
collision-free tail; otherwise the adjoint would have a nonzero
finite-dimensional invariant subspace.  Beyond escape, coefficients branch
with dual norm growing like $C^N$.  Norming functionals therefore produce a
right inverse $R_N$ of $T^N$ with $\|R_N\|\le K C^{-N}<1$.

A contracting right inverse gives shadowing by a convergent correction
series.  Thus $T$ has finite shadowing.  A theorem of Bernardes and
Peris~\cite[Theorem~11]{BP} then upgrades Devaney chaos plus finite shadowing
to \emph{dense} distributional chaos.  The latter is stronger than the
conjectured conclusion.

\section{The theorem}
Let $p'$ denote the conjugate exponent when $p>1$, and put $p'=\infty$ when
$p=1$.

\begin{theorem}\label{thm:main}
Let $1\le p<\infty$, let $f_1,\ldots,f_t:\N\to\N$ be strictly increasing and
pairwise almost disjoint, and let $T$ be given by \eqref{eq:T}.
If $T$ is hypercyclic, then some power $T^N$ has a bounded right inverse
$R_N$ satisfying $\|R_N\|<1$.  Consequently $T$ has the positive and finite
shadowing properties.

If $T$ is Devaney chaotic, then $T$ is densely distributionally chaotic.  In
particular, the conjecture of Bongiorno--Darji--Di Piazza is true.
\end{theorem}

We split the proof into four elementary lemmas.

\begin{lemma}[collision-free tail]\label{lem:tail}
There is $m\in\N$ such that, on
$Y=\overline{\spn}\{e_k:k\ge m\}$, the factor $D$ determined by
\[
                P_YT=DP_Y
\]
has norm
\[
 C:=\|D\|=\|(c_1,\ldots,c_t)\|_{p'}.
\tag{2}\label{eq:C}
\]
If $T$ is hypercyclic, then $C>1$.
\end{lemma}

\begin{proof}
Choose $m$ so that $f_i(k)\ne f_j(k)$ for $i\ne j$ and $k\ge m$.
Strict increase gives $f_i(k)\ge k$, so coordinates $k\ge m$ depend only on
the tail and $P_YT=DP_Y$.  For $k\ge m$, all $f_i(k)$ are distinct.  Moreover,
if $k\ne \ell$, almost disjointness gives
$f_i(k)\ne f_j(\ell)$ for every $i,j$.  Hence all input coordinates appearing
in the tail rows are globally distinct.

Hölder's inequality, followed by this disjointness, gives
\[
 \|Dy\|_p^p
 =\sum_{k\ge m}\left|\sum_i c_i y_{f_i(k)}\right|^p
 \le C^p\sum_{k\ge m}\sum_i|y_{f_i(k)}|^p
 \le C^p\|y\|_p^p.
\]
Equality is attained by putting a dual norming vector for $(c_i)$ on the
children of one row.  The same proof with the maximum norm covers $p=1$.
Thus \eqref{eq:C} holds.

If $x$ is hypercyclic for $T$, then $P_Yx$ is hypercyclic for $D$, since
$P_YT^nx=D^nP_Yx$ and the projection of a dense orbit is dense.  An operator
of norm at most one has a bounded orbit from every vector and cannot be
hypercyclic on the unbounded space $Y$.  Therefore $C>1$.
\end{proof}

For a word $\sigma=(i_1,\ldots,i_r)$, write
$f_\sigma=f_{i_r}\circ\cdots\circ f_{i_1}$ and
$c_\sigma=c_{i_1}\cdots c_{i_r}$.

\begin{lemma}[finite-prefix escape]\label{lem:escape}
Assume that $T$ is hypercyclic.  For every $k<m$ there are $r_k\ge0$, an
index $a_k\ge m$, and a nonzero scalar $\beta_k$ such that $\beta_k$ is the
coefficient of $e_{a_k}^*$ in $(T^*)^{r_k}e_k^*$.
\end{lemma}

\begin{proof}
If no such coefficient existed, the cyclic subspace
\[
 M_k=\spn\{(T^*)^ne_k^*:n\ge0\}
\]
would be a nonzero $T^*$-invariant subspace of
$\spn\{e_1^*,\ldots,e_{m-1}^*\}$, hence finite-dimensional.  Its annihilator
would give a nonzero finite-dimensional factor of $T$.  A dense orbit of $T$
would project to a dense orbit on that factor, contradicting the standard
fact that no linear operator on a nonzero finite-dimensional real or complex
space is hypercyclic.
\end{proof}

For $k\ge m$, set $r_k=0$, $a_k=k$, and $\beta_k=1$.

\begin{lemma}[contracting right inverse of a power]\label{lem:right}
If $T$ is hypercyclic, there are $K<\infty$ and, for every sufficiently large
$N$, a bounded operator $R_N$ such that
\[
                  T^NR_N=I,
       \qquad \|R_N\|\le K C^{-N}.
\tag{3}\label{eq:right}
\]
In particular, $\|R_N\|<1$ for some $N$.
\end{lemma}

\begin{proof}
Fix $N\ge\max_{k<m}r_k$.  Once the path from $k$ reaches $a_k\ge m$, distinct
continuation words never collide.  Indeed, equality of two equal-depth
descendants can be pulled back one level using injectivity and the
cross-input clause in almost disjointness; at a common input at least $m$,
the choice of branch is also unique.  Likewise, equal-depth descendant sets
from distinct starting rows are disjoint.

It follows that the coordinates
\[
      J_{k,N}=\{f_\sigma(a_k):|\sigma|=N-r_k\}
\]
occur in row $k$ of $T^N$ with coefficients
$(\beta_kc_\sigma)_\sigma$, and occur in no other row.  Their $\ell_{p'}$
norm is
\[
       |\beta_k|
       \left(\sum_{|\sigma|=N-r_k}|c_\sigma|^{p'}\right)^{1/p'}
       =|\beta_k|C^{N-r_k},
\]
with the evident maximum-norm formula for $p=1$.  By $\ell_p$--$\ell_{p'}$
duality there is a vector $u_{k,N}$ supported on $J_{k,N}$ such that
\[
 T^Nu_{k,N}=e_k,
 \qquad
 \|u_{k,N}\|_p=|\beta_k|^{-1}C^{-(N-r_k)}.
\tag{4}\label{eq:column}
\]
The supports of the vectors $u_{k,N}$ are pairwise disjoint.  Thus the rule
$R_Ne_k=u_{k,N}$ extends boundedly to $\ell_p$, satisfies $T^NR_N=I$, and
has norm equal to the supremum of the column norms in \eqref{eq:column}.
Since only the indices $k<m$ have exceptional $(r_k,\beta_k)$, \eqref{eq:right}
holds with
\[
 K=\max\left\{1,\max_{k<m}|\beta_k|^{-1}C^{r_k}\right\}.
\]
Now $C>1$ by Lemma~\ref{lem:tail}, so $KC^{-N}<1$ for large $N$.
\end{proof}

\begin{lemma}[right inverse implies shadowing]\label{lem:shadow}
If a bounded operator $A$ on a Banach space has a right inverse $R$ with
$\|R\|=\rho<1$, then $A$ has positive shadowing.  If $T^N$ has positive
shadowing, then $T$ has positive shadowing.
\end{lemma}

\begin{proof}
Let $(x_j)_{j\ge0}$ be a positive pseudo-orbit for $A$ and write
$e_j=x_{j+1}-Ax_j$.  The series
\[
                y=x_0+\sum_{j=0}^\infty R^{j+1}e_j
\tag{5}\label{eq:shadowseries}
\]
converges.  Since $AR=I$, cancellation of the first $n$ error terms gives
\[
 A^ny-x_n=\sum_{j=n}^\infty R^{j+1-n}e_j,
 \qquad
 \|A^ny-x_n\|\le \frac{\rho}{1-\rho}\sup_j\|e_j\|.
\]
This proves positive shadowing for $A$.

For the second assertion, group a $\delta$-pseudo-orbit of $T$ into blocks of
length $N$.  Its sampled sequence at times $0,N,2N,\ldots$ is a
$\delta\sum_{r=0}^{N-1}\|T\|^r$-pseudo-orbit for $T^N$.  Shadow the sampled
sequence, then propagate over each of the finitely many intermediate times.
The estimate
\[
 \|T^r(T^{jN}y-x_{jN})\|
 +\delta\sum_{s=0}^{r-1}\|T\|^s,
 \qquad 0\le r<N,
\]
is made arbitrarily small by first choosing the shadowing tolerance for
$T^N$ and then $\delta$.  Hence $T$ has positive shadowing.
\end{proof}

\begin{proof}[Proof of Theorem~\ref{thm:main}]
Lemmas~\ref{lem:tail}--\ref{lem:shadow} show that a hypercyclic $T$ has
positive shadowing.  Positive shadowing immediately implies finite
shadowing.  If $T$ is Devaney chaotic, Bernardes and Peris
\cite[Theorem~11]{BP} apply on the Fréchet space $\ell_p$ and imply that $T$
is densely distributionally chaotic.  This proves, and strengthens, the
source conjecture.
\end{proof}

\section{Verification and limitations}
The exact-arithmetic script \texttt{code/verify\_right\_inverse.py} checks the
dual normalization, disjoint supports, and $T^NR_N=I$ for a finite two-branch
tree with rational coefficients.  It is only a sanity check; the proof above
is symbolic.

The theorem does not assert that every formal sum \eqref{eq:T} is chaotic; it
starts from the conjecture's exact conditional hypothesis.  In particular,
it does not rely on the source paper's separate sufficient condition for
chaos.  The proof also genuinely uses the source's strong cross-input
almost-disjointness condition, not merely eventual pointwise inequality of
the functions.

\paragraph{Novelty check.}
On 12 August 2026, bounded searches of arXiv and the run indexes used the
exact title and arXiv id together with ``distributionally chaotic'',
``shadowing'', ``finite shadowing'', ``pseudo-shift'', and ``generalized
backward shift''.  They found the general Bernardes--Peris theorem but no
paper explicitly resolving this conjecture or giving the contracting-power
right-inverse argument for these almost-disjoint sums.  Novelty confidence is
therefore moderate, not definitive.

\paragraph{Human-review recommendation.}
Prioritize review.  The central points to check are the escape lemma and the
absence of descendant collisions used in Lemma~\ref{lem:right}; after those,
the quantitative contraction and shadowing steps are routine.

\begin{thebibliography}{9}
\bibitem{BDD}
D.~Bongiorno, U.~B. Darji, and L.~Di Piazza,
\emph{Rolewicz-type chaotic operators}, arXiv:1504.02445 (2015).

\bibitem{BP}
N.~C. Bernardes Jr. and A.~Peris,
\emph{On shadowing and chain recurrence in linear dynamics},
arXiv:2305.02714 (2023), Theorem~11.
\end{thebibliography}

\end{document}
